\documentclass[11pt]{article}
\usepackage[margin=1in]{geometry}
\usepackage{amsmath,amssymb,amsthm}
\usepackage[hidelinks]{hyperref}
\newtheorem{lemma}{Lemma}
\newtheorem{theorem}[lemma]{Theorem}
\newtheorem{corollary}[lemma]{Corollary}
\newtheorem{proposition}[lemma]{Proposition}
\begin{document}
\setlength{\emergencystretch}{1em}
\title{Boundary-degree Frobenius descent and the ordinary support $\{2,6\}$}
\author{}\date{}\maketitle
\begin{abstract}
We extend the controlled-radius Frobenius noncollision invariant to
the missing boundary degree, with a precise constant-coefficient
valuation condition. The extension applies to a mixed-adjoint relation
for lower exponent two and fourth-power gap. It proves consecutive
coprimality at every ordinary order for the support $\{2,6\}$.
The same proof gives an explicit sufficient index condition at
prime-power gaps. It does not establish unrestricted binomial SFC.
\end{abstract}

\section{The exact stage class and descent input}
Fix a prime $p$ and a complete algebraically closed extension $K$ of
$\mathbb Q_p$, with valuation $v(p)=1$, ring $\mathcal O$, maximal
ideal $\mathfrak m$, and residue field $k$. Put
\[
 q_s=p^s,\qquad b_s=\frac{p^s-1}{p-1},\qquad
 B_s=b_s-s,\qquad
 \Pi_{n,q}(t)=\prod_{\ell=0}^{n-1}(t+q\ell).
\]
Here $b_s\ge s$, so $B_s\ge0$.
At stage $s\ge1$ use precisely the closed class
\[
 A=c-\theta+F(z),\quad\theta=zD_z,\qquad
 c\in\mathbb Z_p+p^s\mathcal O,
\]
\[
 F(z)=z+\sum_{r\ge1}a_rz^{p^r},\quad a_r\in p\mathbb Z_p,\qquad
 h\in\mathcal O[[z]],\quad
 \bar h=\overline{h(0)}\in k^\times.
\]
The series and their finite Euler derivatives converge at $v(z)>0$.
For $f(z)=\sum f_jz^j$ its cyclic column is
\[
 \mathcal C_q(f;u)_j=\sum_{\ell\ge0}f_{q\ell+j}u^\ell,
             \qquad0\le j<q.
\]
When $u\ne0$ and $v(u)>0$, its vanishing is equivalent to
$f(z)=0$ at all $q$ roots of $z^q=u$.
Multiplying by a convergent unit series preserves this vanishing.
All operators below act on whole series before evaluation.
No characteristic-zero operator action on a fixed quotient
by $z^q-u$ is assumed.

We import the exact Artin--Hasse stage change of Lemma
\texttt{lem:artin-hasse-closure} in the
\href{ARITHMETIC_FRONTIER_TERMINAL_ORDERS.tex}{arithmetic note}.
Here is its complete interface. The integral unit gauge
\[
 E(z)=\exp\left(\sum_{r\ge0}\frac{z^{p^r}}{p^r}\right)
\]
has integral inverse and conjugates $A$ to
$c-\theta+\sum_{r\ge1}(a_r-1)z^{p^r}$.
For $s\ge2$, write $q=pq'$, $q'=q_{s-1}$,
$c=c_0+p^s\eta$, $c_0\in\mathbb Z_p$, $\eta\in\mathcal O$.
For residue $j$ modulo $p$, choose $e_j\in\{0,\ldots,p-1\}$
congruent to $c_0-j$ modulo $p$. Put
\[
 b=a_1-1\in\mathbb Z_p^\times,\quad
 c_j=(c-j-e_j)/p,\quad
 w=bz^p/p,\quad u_*=(b/p)^{q'}u.
\]
The block operator is $e_j+pL_j$, where
\[
 L_j=c_j-\theta_w+w+
 \sum_{r\ge1}\frac{(a_{r+1}-1)p^{p^r-1}}{b^{p^r}}w^{p^r}.
\]
It belongs to the same class at stage $s-1$.
If $E^{-1}h=\sum_{j=0}^{p-1}z^j\widetilde h_j(z^p)$,
its new seed is $h_j^*(w)=\widetilde h_j(pw/b)$ and
\[
 \overline{h_j^*}
 =\overline{h(0)}\,(-1)^j/j!\in k^\times,\qquad
 v(u_*)\ge b_s-q'=b_{s-1}>0.
\]
Consequently a zero cyclic column at stage $s$ gives, for each block,
\[
 \mathcal C_{q'}\!\left(
 \Pi_{n,q}(e_j+pL_j)P(e_j+pL_j)h_j^*;u_*\right)=0.       \tag{1}\label{eq:block}
\]
These integrality, seed, radius and block identities are the only
descent facts imported here; the boundary-degree induction follows.

\section{The boundary-degree invariant}

\begin{theorem}[Boundary-degree noncollision]\label{thm:boundary}
Let $s,n\ge1$, $u\ne0$, and $v(u)\ge b_s$, in the stage class above.
Suppose $P\in\mathcal O[t]$ is primitive and
\[
 \deg\bar P\le q_s-1,\qquad v(P(0))>B_s=b_s-s.
\]
Then
\[
 \mathcal C_{q_s}\bigl(\Pi_{n,q_s}(A)P(A)h;u\bigr)\ne0.
\]
The ordinary degree of $P$ is unrestricted.
The convention $v(0)=+\infty$ is allowed at its constant coefficient.
\end{theorem}
\begin{proof}
We induct on $s$. At $s=1$, reduction has $z^p=0$,
$\bar A=z-\theta+\bar c$, $\bar c\in\mathbb F_p$,
and $\Pi_{n,p}(\bar A)=\bar A^n$.
The seed is a constant unit.
The triangular diagonal of $\bar A$ is $\bar c-j$ for $0\le j<p$,
so its characteristic polynomial is $X^p-X$.
Moreover $1,\bar A1,\ldots,\bar A^{p-1}1$ form a basis,
being monic in degrees zero through $p-1$.
Thus $1$ has a nonzero component in every simple eigenspace,
including each nonzero eigenvalue in $\mathbb F_p^\times$.
A collision would force $\bar P$ to vanish at all these $p-1$ points.
Primitivity and $\deg\bar P\le p-1$ then force
$\bar P=a(t^{p-1}-1)$ with $a\in k^\times$.
Its constant coefficient is nonzero, contradicting
$v(P(0))>B_1=0$.

Let $s\ge2$ and suppose a collision occurs.
Use \eqref{eq:block}. In a nonresonant block $e_j\ne0$, absorb
$\Pi_{n,q}(e_j+pL_j)$ into its seed $\widehat h_j$.
This seed is integral and reduces to the constant unit
$e_j^n\overline{h_j^*}$.
The $q'$ vectors
\[
 v_r=\mathcal C_{q'}(L_j^r\widehat h_j;u_*),
                  \qquad0\le r<q',
\]
form an integral unit basis. Indeed reduction sets $u_*=0$
and the seed to a constant, and
$(w-\theta+\overline{c_j})^r1$ is monic of degree $r<q'$.
All subsequent vectors $v_r$ are integral as well.
With $P^{[r]}(e)=[t^r]P(e+t)$, the finite divided-Taylor
expansion gives
\[
 \sum_{r\ge0}p^rP^{[r]}(e_j)v_r=0.
\]
This is finite even when the ordinary degree of $P$ is large.
Terms with $r\ge q'$ have valuation at least $q'$.
Inverting the unit basis therefore gives
\[
 v(P^{[r]}(e_j))\ge q'-r>0,\qquad0\le r<q'.
\]
As $j$ runs through the nonresonant blocks, $e_j$ runs through
$\mathbb F_p^\times$. Hence
\[
 \bar P(t)=(t^{p-1}-1)^{q'}R(t),\qquad
                 0\ne R,\quad\deg R\le q'-1.            \tag{2}\label{eq:residualfactor}
\]
The prefactor has nonzero constant coefficient.
Thus the index $r_0$ of the first unit coefficient of $P$
is at most $q'-1$.

Write $P=\sum c_rt^r$ and choose an actual nonzero coefficient
\[
 \lambda=p^{r_*}c_{r_*},\qquad
 M=v(\lambda)=\min_r\{v(c_r)+r\}\le r_0\le q'-1.
\]
The minimum may be fractional. We divide by this actual field element,
not by an unspecified symbol $p^M$.
Then $Q(t)=P(pt)/\lambda$ is integral and primitive.
For $r\ge q'$,
\[
 v([t^r]Q)=v(c_r)+r-M\ge r-M\ge1.
\]
Consequently $\deg\bar Q\le q'-1$. Also
\[
 v(Q(0))=v(P(0))-M
 >B_s-(q'-1)=B_{s-1}.                                  \tag{3}\label{eq:constant-transfer}
\]
There is exactly one resonant block $e_j=0$.
Here $\Pi_{n,q}(pL_j)=p^n\Pi_{n,q'}(L_j)$.
After dividing \eqref{eq:block} by $p^n\lambda$, we obtain
\[
 \mathcal C_{q'}\bigl(
       \Pi_{n,q'}(L_j)Q(L_j)h_j^*;u_*\bigr)=0.
\]
All stage, seed and radius conditions are unchanged, and
\eqref{eq:constant-transfer} is the boundary condition at stage $s-1$.
This contradicts induction.
\end{proof}

\begin{proposition}[The strict bound is sharp at every stage]
\label{prop:sharpness}
For every prime $p$ and $s\ge1$, there is a stage-$s$ example
with $n=1$, primitive $P\in\mathbb Z[t]$ of degree $q_s-1$,
and $v(P(0))=B_s$, for which the cyclic column in
Theorem~\ref{thm:boundary} vanishes.
Thus its strict constant-coefficient inequality cannot be replaced
by $v(P(0))\ge B_s$ in the same class.
\end{proposition}
\begin{proof}
Put $q=p^s$, $A=q-\theta+z$, $h=1$, and
\[
 u=(-1)^q q!,\qquad
 \chi(t)=\prod_{j=1}^{q}(t-j)-(-1)^q q!=tP(t).
\]
Then $P$ is monic with integer coefficients and degree $q-1$.
Since
\[
 (A-(q-m))z^m=z^{m+1}\qquad(0\le m<q),
\]
successive application, in the commuting product, gives
$\prod_{j=1}^{q}(A-j)1=z^q$.
Consequently
\[
 \Pi_{1,q}(A)P(A)1=\chi(A)1=z^q-u,
\]
whose cyclic column at $u$ is zero.
Legendre's formula gives $v_p(q!)=b_s$, so the radius condition
is exact. With $H_q=\sum_{j=1}^{q}1/j$,
\[
 P(0)=\chi'(0)=(-1)^{q-1}q!H_q.
\]
The term $1/q$ is the unique term of least valuation in $H_q$,
so $v_p(H_q)=-s$ and $v_p(P(0))=b_s-s=B_s$.
The operator and seed belong to the stated stage class.
This is a sharpness example for the noncollision invariant,
not an SFC counterexample.
\end{proof}

\section{Application to lower exponent two}

The ordinary factorial functional is $\mathcal L_1(X^j)=j!$.
For an integer gap $q\ge3$ put
\[
 P_{j,q}(t)=\sum_{k=0}^j\binom jk(2j+1)_{qk}t^k .
\]
For fixed lower order $n\ge1$ set
\[
 \gamma=2n+1,\qquad
 \kappa_q=\frac{q}{(q+2)((q+2)n+q+1)},
\]
and use the auxiliary recurrence
\[
 B_{0,\gamma,q}=1,\qquad
 B_{j+1,\gamma,q}
  =(z+qj+\gamma+q-1)B_{j,\gamma,q}
       -zB_{j,\gamma,q}' .
\]
The precise input is Theorem \texttt{thm:mixed} of the
\href{LOWER_TWO_PRIME_GAP_SFC.tex}{general mixed-rank note}.
Its necessary direction states:
if $P_{n,q}(t)=P_{n+1,q}(t)=0$, then for
$u=(-1)^{q+1}/t$ the $q$-cyclic columns of
\[
 B_{n+1,\gamma,q},\ldots,B_{n+q-2,\gamma,q},\quad
 B_{n+q-1,\gamma,q}+\kappa_q z^q B_{n,\gamma,q}            \tag{4}\label{eq:mixed-input}
\]
have rank less than $q-1$.
Only this algebraic mixed-rank identity, valid in characteristic zero
and at multiple lower roots, is used.

\begin{corollary}[An explicit prime-power index condition]
\label{cor:prime-power-index}
Let $q=p^s\ge3$, $n\ge1$, and put $b_s=(q-1)/(p-1)$.
If
\[
 v_p(n)>b_s-2s,\qquad v_p(\kappa_q)>-s,                   \tag{5}\label{eq:indexcondition}
\]
then $P_{n,q},P_{n+1,q}$ are coprime over $\mathbb Q$ and
over $\mathbb C$.
\end{corollary}
\begin{proof}
Suppose there is a common root over $K$. It is nonzero.
Choose $z^q=u=(-1)^{q+1}/t$.
Every interval of $qk$ consecutive integers contains exactly
$qk/p^j$ multiples of $p^j$ for $1\le j\le s$, so
\[
 v_p((2n+1)_{qk})\ge b_sk.
\]
The monic reciprocal polynomial in $z$ therefore gives
$v(u)\ge b_s$: below that radius its leading term is uniquely least.

The rank relation \eqref{eq:mixed-input} supplies coefficients
$a_1,\ldots,a_{q-1}\in\mathcal O$, at least one a unit, for a zero
linear combination of the columns. Put
\[
 A=\gamma+q-1-zD_z+z,\qquad
 R_i(t)=\prod_{\ell=n}^{n+i-1}(t+q\ell).
\]
The stage class holds with $c=\gamma+q-1\in\mathbb Z_p$,
all higher coefficients zero, and seed $h=1$.
The recurrence gives
\[
 B_{n+i,\gamma,q}=\Pi_{n,q}(A)R_i(A)1.
\]
Consequently the zero linear combination is the cyclic column of
$\Pi_{n,q}(A)P(A)1$, where
\[
 P(t)=\sum_{i=1}^{q-1}a_iR_i(t)+a_{q-1}\kappa_q u.       \tag{6}\label{eq:relation-polynomial}
\]
Here $u$ in the constant coefficient is the fixed evaluated field
element. It is not a variable on which the Euler operator acts.

The hypotheses imply
$v(\kappa_q u)>b_s-s=B_s\ge0$.
Also $R_i\bmod\mathfrak m=t^i$, so
$\bar P=\sum_{i=1}^{q-1}\bar a_it^i\ne0$:
$P$ is primitive with $\deg\bar P\le q-1$.
Its constant coefficient satisfies
\[
 v(P(0))\ge
 \min\{s+v_p(n),\,v_p(\kappa_q)+b_s\}>B_s,
\]
because $R_i(0)=q^i(n)_i$ has valuation at least $s+v_p(n)$.
Theorem~\ref{thm:boundary} contradicts the zero cyclic column.
Thus the two rational polynomials have no common root over $K$;
a rational Bezout identity gives coprimality over $\mathbb C$.
\end{proof}

\begin{theorem}[The ordinary support $\{2,6\}$]\label{thm:fourth-gap}
Every nonzero complex phase $X^2(A+BX^4)$ satisfies Strong Factorial:
its ordinary factorial moments cannot both vanish at any consecutive
positive orders.
\end{theorem}
\begin{proof}
Take $p=2,s=2,q=4$. Then $b_s=3$, $B_s=1$, and
\[
 \kappa_4=\frac2{3(6n+5)},\qquad v_2(\kappa_4)=1.
\]
For every integer $n\ge1$, one has $v_2(n)\ge0>3-4=-1$
and $v_2(\kappa_4)=1>-2$.
Corollary~\ref{cor:prime-power-index} applies at every order.

For $AB\ne0$, the exact normalization is
\[
 \mathcal L_1((X^2(A+BX^4))^j)
       =A^j(2j)!P_{j,4}(B/A).
\]
If $A=0$ or $B=0$, every positive moment is nonzero because
the phase is a nonzero monomial.
\end{proof}

For instance, \eqref{eq:indexcondition} also gives a uniform
sufficient index family at gap $q=2^s$, $s\ge3$:
\[
                    2^{\,2^s-2s}\mid n.
\]
Indeed $v_2(\kappa_q)=s-1$ there. These are sufficient subfamilies,
not an all-order theorem at every prime-power gap.

\paragraph{Replay and scope.}
The checker replays the boundary algebra, exact constant transfer,
literal relation-polynomial controls, and the fourth-gap valuation
conditions. Universal noncollision follows from the induction,
not a finite parameter search.
This note proves support $\{2,6\}$ and the stated sufficient index
criterion. It does not prove all ordinary binomials of degree six,
unrestricted binomial SFC, general FC/SFC, or every Gamma shape.
Internal replay is distinct from formal verification, external
review and a priority claim.
\end{document}
